%----------------------------------------------------------
% 제10장. 모델 개발 단계의 거버넌스
%----------------------------------------------------------

\chapter{모델 개발 단계의 거버넌스}
\label{ch:model_development}

모델 개발 단계는 데이터 사이언티스트와 엔지니어들이 가장 활발하게 활동하는 영역이자, AI 시스템의 성능과 안전성이 결정되는 시기이다. 이 단계에서의 거버넌스는 ``창의성을 저해하는 통제''가 아니라, ``시행착오를 줄이고 성공 확률을 높이는 가이드라인''으로 작동해야 한다. 모델이 복잡해질수록(블랙박스(Black-box)), 개발 과정의 투명성과 검증의 중요성은 더욱 커진다.

본 장에서는 모델 개발의 전 과정---표준화된 문서화, 실험 관리와 재현성\cite{studer2021towards}, 공정성 검증, 강건성\cite{madry2018towards} 및 보안 테스트, 설명가능성 구현, 승인과 릴리스 관리, 그리고 기술 부채 관리---에 걸쳐 실무에 바로 적용 가능한 거버넌스 프레임워크를 제시한다.

%----------------------------------------------------------
\section{모델 개발 표준과 문서화}
\label{sec:10.1}
%----------------------------------------------------------

개발자마다 제각각의 방식으로 코드를 작성하고 모델을 만든다면, 그 모델은 ``일회용''이 되고 만다. 표준화는 협업과 유지보수의 기반\cite{paleyes2022challenges}이다.

\subsection{모델 카드(Model Card) 작성}
\label{sec:10.1.1}

구글이 제안한 모델 카드\cite{mitchell2019model}는 AI 모델의 ``사양서''이자 ``설명서''\cite{arnold2019factsheets}이다. 모델 개발 완료 시 다음 내용을 필수적으로 포함하여 문서화해야 한다. 이는 향후 감사의 핵심 증거 자료가 된다.

\begin{enumerate}
  \item \textbf{모델 세부 정보(Model Details)}:
  \begin{itemize}
    \item 개발자/조직, 모델 버전, 모델 아키텍처(예: 트랜스포머(Transformer), CNN), 파라미터 수.
    \item 모델 생성 일자, 라이선스 정보.
  \end{itemize}
  \item \textbf{의도된 용도(Intended Use)}:
  \begin{itemize}
    \item 모델이 해결하고자 하는 문제\cite{rismani2023beyond}와 타겟 사용자.
    \item 사용하면 안 되는 경우(Out-of-scope Use) 명시 (예: ``이 챗봇은 의학적 조언을 제공하지 않음'').
  \end{itemize}
  \item \textbf{요인(Factors)}:
  \begin{itemize}
    \item 모델 성능에 영향을 미치는 변수들 (인구통계학적 그룹, 환경 조건, 입력 기기 등).
  \end{itemize}
  \item \textbf{평가 지표(Metrics)}:
  \begin{itemize}
    \item 성능 측정에 사용된 지표(정확도(Accuracy), F1-점수, AUC 등)와 그 결과값.
    \item 성능의 신뢰 구간(Confidence Interval).
  \end{itemize}
  \item \textbf{학습 데이터(Training Data)}:
  \begin{itemize}
    \item 사용된 데이터셋의 출처, 크기, 전처리 방식.
  \end{itemize}
  \item \textbf{윤리적 고려사항(Ethical Considerations)}:
  \begin{itemize}
    \item 식별된 리스크(편향성 등)와 이를 완화하기 위해 취한 조치.
  \end{itemize}
\end{enumerate}

\subsection{모델 카드 작성 자동화}
\label{sec:10.1.2}

모델 카드를 수작업으로 작성하면 누락과 불일치가 빈번하게 발생한다. 따라서 학습 파이프라인과 연동하여 모델 카드를 자동 생성하는 체계가 필요하다. 다음 코드는 모델 학습 완료 시 자동으로 모델 카드를 생성하는 시스템을 보여준다.

\begin{lstlisting}[language=Python, caption={모델 카드 자동 생성기}, label={lst:model_card_auto}]
import json
import datetime
from dataclasses import dataclass, field, asdict
from typing import Dict, List, Optional

@dataclass
class ModelCardGenerator:
    """모델 카드 자동 생성 시스템"""
    model_name: str
    model_version: str
    developer: str
    model_type: str
    description: str
    intended_use: str
    out_of_scope_use: str
    training_data_source: str
    training_data_size: int
    metrics: Dict[str, float] = field(default_factory=dict)
    fairness_metrics: Dict[str, float] = field(
        default_factory=dict
    )
    ethical_considerations: List[str] = field(
        default_factory=list
    )
    limitations: List[str] = field(default_factory=list)
    creation_date: str = field(
        default_factory=lambda: datetime.date.today().isoformat()
    )

    def add_performance_metrics(
        self, y_true, y_pred, y_prob=None
    ):
        """학습 결과에서 성능 지표를 자동 추출"""
        from sklearn.metrics import (
            accuracy_score, precision_score,
            recall_score, f1_score, roc_auc_score
        )
        self.metrics["accuracy"] = round(
            accuracy_score(y_true, y_pred), 4
        )
        self.metrics["precision"] = round(
            precision_score(y_true, y_pred, average="weighted"), 4
        )
        self.metrics["recall"] = round(
            recall_score(y_true, y_pred, average="weighted"), 4
        )
        self.metrics["f1_score"] = round(
            f1_score(y_true, y_pred, average="weighted"), 4
        )
        if y_prob is not None:
            self.metrics["auc"] = round(
                roc_auc_score(y_true, y_prob), 4
            )

    def add_fairness_evaluation(
        self, y_true, y_pred, sensitive_attr
    ):
        """민감 속성 기반 공정성 지표\cite{mehrabi2021survey} 자동 계산"""
        import numpy as np
        groups = np.unique(sensitive_attr)
        group_rates = {}
        for g in groups:
            mask = sensitive_attr == g
            group_rates[str(g)] = np.mean(y_pred[mask])
        max_rate = max(group_rates.values())
        min_rate = min(group_rates.values())
        self.fairness_metrics["disparate_impact"] = round(
            min_rate / max_rate if max_rate > 0 else 0, 4
        )
        self.fairness_metrics["group_rate_diff"] = round(
            max_rate - min_rate, 4
        )
        self.fairness_metrics["group_rates"] = group_rates

    def generate(self, output_path: str = None) -> dict:
        """모델 카드를 JSON 형식으로 생성"""
        card = {
            "schema_version": "1.0",
            "model_details": {
                "name": self.model_name,
                "version": self.model_version,
                "developer": self.developer,
                "type": self.model_type,
                "description": self.description,
                "creation_date": self.creation_date,
            },
            "intended_use": {
                "primary_use": self.intended_use,
                "out_of_scope": self.out_of_scope_use,
            },
            "training_data": {
                "source": self.training_data_source,
                "size": self.training_data_size,
            },
            "performance": self.metrics,
            "fairness": self.fairness_metrics,
            "ethical_considerations": self.ethical_considerations,
            "limitations": self.limitations,
        }
        if output_path:
            with open(output_path, "w", encoding="utf-8") as f:
                json.dump(card, f, ensure_ascii=False, indent=2)
        return card

# 사용 예시
generator = ModelCardGenerator(
    model_name="CreditScoring-v2",
    model_version="2.1.0",
    developer="리스크모델링팀 김이박",
    model_type="XGBoost Classifier",
    description="소상공인 신용 대출 심사 보조 모델",
    intended_use="소상공인 대출 신청의 신용 위험 평가",
    out_of_scope_use="개인 가계 대출 심사에 사용 불가",
    training_data_source="2024-2025 대출 이력 데이터",
    training_data_size=100000,
)
card = generator.generate("model_card_credit_v2.json")
print(f"모델 카드 생성 완료: {card['model_details']['name']}")
\end{lstlisting}

\subsection{데이터시트(Datasheets for Datasets) 작성 요구사항}
\label{sec:10.1.3}

\cite{gebru2021datasheets}가 제안한 데이터시트는 모델 카드의 ``데이터 버전''이라 할 수 있다. 데이터셋의 생성 동기, 수집 과정, 전처리 방법, 배포 조건 등을 체계적으로 문서화함으로써 데이터 투명성을 확보한다. 다음은 데이터시트에 반드시 포함해야 할 핵심 항목들이다.

\begin{table}[htbp]
\centering
\caption{데이터시트 필수 작성 항목}
\label{tab:datasheet_requirements}
\begin{tabularx}{\textwidth}{p{3cm}X}
\toprule
\textbf{항목} & \textbf{상세 내용} \\
\midrule
동기(Motivation) & 데이터셋이 생성된 목적, 자금 지원 주체, 기존 데이터셋과의 차별점 \\
\midrule
구성(Composition) & 인스턴스 수, 데이터 유형(텍스트/이미지/정형), 레이블 존재 여부, 민감 정보 포함 여부, 결측값 비율 \\
\midrule
수집 과정(Collection) & 수집 메커니즘(API, 크롤링, 설문 등), 수집 기간, 수집 대상의 동의 획득 방식 \\
\midrule
전처리(Preprocessing) & 적용된 전처리 기법, 원본 데이터 보존 여부, 익명화 처리 방식 \\
\midrule
용도(Uses) & 권장 사용 용도, 부적절한 사용 사례, 향후 사용 시 주의사항 \\
\midrule
배포(Distribution) & 배포 라이선스, 접근 제한 조건, 저작권 정보 \\
\midrule
유지보수(Maintenance) & 관리 주체, 업데이트 주기, 오류 보고 채널 \\
\bottomrule
\end{tabularx}
\end{table}

데이터시트 작성은 모델 개발 시작 전에 완료되어야 하며, 데이터 변경 시마다 갱신해야 한다. 특히 민감 정보 포함 여부, 동의 획득 방식, 대표성 분석 결과는 법무 검토 대상이 된다.

\subsection{개발 환경 표준화}
\label{sec:10.1.4}

모델 개발 환경의 비표준화는 ``내 컴퓨터에서는 됩니다'' 문제의 원인이 된다. 거버넌스 관점에서 개발 환경 표준화는 재현성과 보안의 기반이다.

\paragraph{컨테이너화(Containerization)}
도커(Docker) 기반 컨테이너화는 개발 환경 표준화의 핵심이다. 모든 모델 개발 프로젝트는 다음을 포함하는 표준 도커 이미지를 기반으로 시작해야 한다.

\begin{itemize}
  \item 승인된 베이스 이미지 (보안팀 검증 완료)
  \item 고정된 파이썬(Python) 버전 및 CUDA 버전
  \item 의존성 패키지의 정확한 버전 명시 (\texttt{requirements.txt} 또는 \texttt{poetry.lock})
  \item 보안 취약점 스캔 통과 확인
\end{itemize}

\paragraph{의존성 관리(Dependency Management)}
ML 프로젝트의 의존성은 일반 소프트웨어보다 복잡하다. GPU 드라이버, CUDA 툴킷, 딥러닝 프레임워크 간의 호환성 매트릭스를 관리해야 한다.

\begin{itemize}
  \item \texttt{pip freeze} 또는 \texttt{conda list --export}를 통한 전체 의존성 스냅샷 저장
  \item 내부 패키지 레지스트리(예: Artifactory, Nexus)를 통한 승인된 패키지만 사용
  \item 주기적인 의존성 보안 취약점 스캔 (예: \texttt{pip-audit}, \texttt{safety check})
\end{itemize}

\subsection{코드 리뷰 및 피어 리뷰 체계}
\label{sec:10.1.5}

ML 코드는 일반 소프트웨어 코드와 다른 리뷰 관점이 필요하다. 데이터 전처리 로직의 올바름, 평가 지표 계산의 정확성, 데이터 누수(Data Leakage) 여부 등 ML 고유의 검증 포인트가 존재한다.

\begin{table}[htbp]
\centering
\caption{ML 코드 리뷰 체크리스트}
\label{tab:ml_code_review}
\begin{tabularx}{\textwidth}{p{3cm}X}
\toprule
\textbf{검증 영역} & \textbf{체크 포인트} \\
\midrule
데이터 누수 & 학습 데이터와 테스트 데이터 간 정보 유출 없는지 확인. 시계열 데이터의 경우 미래 정보 사용 여부 확인 \\
\midrule
전처리 일관성 & 학습 시 전처리와 추론 시 전처리가 동일한 코드를 사용하는지 확인 \\
\midrule
재현성 & 랜덤 시드 고정 여부, 데이터 셔플링(Shuffling) 시드 일관성 \\
\midrule
하드코딩 & 데이터 경로, 하이퍼파라미터 등이 설정 파일로 분리되어 있는지 확인 \\
\midrule
평가 정확성 & 성능 지표 계산 로직이 올바른지, 클래스 불균형 고려 여부 \\
\midrule
리소스 관리 & GPU 메모리 해제, 대용량 데이터 배치 처리, 메모리 누수 여부 \\
\bottomrule
\end{tabularx}
\end{table}

피어 리뷰는 최소 2인 이상이 참여해야 하며, 리뷰어 중 최소 1인은 해당 도메인 전문가여야 한다. 고위험 등급 모델의 경우 독립적인 제3자 리뷰를 추가로 수행한다.

%----------------------------------------------------------
\section{실험 관리와 재현성}
\label{sec:10.2}
%----------------------------------------------------------

ML 모델 개발은 본질적으로 실험의 연속이다. 수십, 수백 번의 실험을 체계적으로 관리하지 않으면 ``어떤 설정이 최적이었는지'' 알 수 없게 된다. 재현성(Reproducibility)은 과학적 엄밀성의 기본이자, 규제 감사 대응의 핵심이다.

\subsection{실험 추적 체계}
\label{sec:10.2.1}

실험 추적은 모든 실험의 입력(데이터, 하이퍼파라미터, 코드 버전), 과정(학습 로그), 출력(성능 지표, 모델 아티팩트)을 기록하고 비교할 수 있게 한다. 대표적인 도구로 MLflow와 W\&B(Weights \& Biases)가 있다.

\paragraph{MLflow 연계} MLflow는 오픈소스 실험 관리 플랫폼으로, 실험 추적(Tracking), 모델 패키징(Projects), 모델 레지스트리(Registry) 기능을 제공한다. 거버넌스 관점에서 다음이 중요하다.
\begin{itemize}
  \item 모든 실험은 고유한 실행 ID(Run ID)로 추적
  \item 하이퍼파라미터, 지표, 아티팩트가 자동으로 로깅
  \item 모델 레지스트리를 통한 모델 버전 관리와 스테이지 전환(Staging → Production)
  \item 태그(Tag)를 활용한 거버넌스 메타데이터 기록 (리뷰 상태, 승인자 등)
\end{itemize}

\paragraph{W\&B(Weights \& Biases) 연계} W\&B는 실험 시각화와 협업에 강점이 있다. 특히 다음 기능이 거버넌스에 유용하다.
\begin{itemize}
  \item 실시간 실험 대시보드 공유로 팀 전체의 가시성 확보
  \item 데이터셋 버전 관리(Artifacts)와 계보(Lineage) 추적
  \item 스윕(Sweeps)을 통한 하이퍼파라미터 탐색 기록
  \item 리포트(Reports) 기능을 통한 실험 결과 문서화
\end{itemize}

\subsection{하이퍼파라미터 관리와 버전 관리}
\label{sec:10.2.2}

하이퍼파라미터(Hyperparameter)는 모델 성능에 직접적인 영향을 미치는 설정값이다. 이를 코드에 하드코딩하는 것은 재현성과 추적성을 모두 해치는 안티패턴(Anti-pattern)이다.

\begin{itemize}
  \item \textbf{설정 파일 분리}: YAML 또는 JSON 형식의 설정 파일로 하이퍼파라미터를 분리하고, 버전 관리 시스템(Git)으로 관리한다.
  \item \textbf{설정 스키마 검증}: 설정 파일의 유효성을 자동 검증하여 잘못된 값 입력을 방지한다.
  \item \textbf{탐색 이력 기록}: 그리드 서치(Grid Search), 베이지안 최적화(Bayesian Optimization) 등 하이퍼파라미터 탐색 과정의 전체 이력을 기록한다.
\end{itemize}

\subsection{재현성 보장을 위한 체크리스트}
\label{sec:10.2.3}

재현성은 ``동일한 입력과 설정으로 동일한 결과를 얻을 수 있는 능력''이다. ML 실험의 재현성을 보장하기 위해 다음 체크리스트를 적용한다.

\begin{table}[htbp]
\centering
\caption{재현성 보장 체크리스트}
\label{tab:reproducibility_checklist}
\begin{tabularx}{\textwidth}{p{0.5cm}p{3.5cm}X}
\toprule
 & \textbf{항목} & \textbf{상세 요구사항} \\
\midrule
1 & 랜덤 시드 고정 & Python, NumPy, PyTorch/TensorFlow, CUDA의 모든 랜덤 시드를 고정하고 기록 \\
\midrule
2 & 환경 스냅샷 & 실행 시점의 전체 패키지 버전, OS 정보, GPU 드라이버 버전 기록 \\
\midrule
3 & 데이터 버전 관리 & DVC(Data Version Control) 등을 활용하여 데이터셋의 정확한 버전을 추적 \\
\midrule
4 & 코드 버전 관리 & 실험 실행 시점의 Git 커밋 해시를 기록 \\
\midrule
5 & 하이퍼파라미터 기록 & 모든 하이퍼파라미터를 설정 파일로 관리하고 실험에 연결 \\
\midrule
6 & 하드웨어 정보 & GPU 모델, CPU 코어 수, 메모리 크기 등 하드웨어 스펙 기록 \\
\midrule
7 & 비결정론적 연산 제어 & cuDNN 비결정론적 알고리즘 비활성화 여부 기록 \\
\bottomrule
\end{tabularx}
\end{table}

\subsection{재현성 보장 실험 관리자}
\label{sec:10.2.4}

다음 코드는 위 체크리스트의 핵심 항목들을 자동으로 적용하고 기록하는 실험 관리자 클래스를 보여준다.

\begin{lstlisting}[language=Python, caption={재현성 보장 실험 관리자}, label={lst:reproducibility_manager}]
import os
import random
import hashlib
import json
import subprocess
import platform
import datetime
from typing import Dict, Any, Optional

class ReproducibilityManager:
    """ML 실험의 재현성을 보장하는 관리자"""

    def __init__(self, seed: int = 42):
        self.seed = seed
        self.env_snapshot = {}
        self.experiment_log = {}

    def fix_all_seeds(self):
        """모든 랜덤 시드를 고정"""
        import numpy as np
        random.seed(self.seed)
        np.random.seed(self.seed)
        os.environ["PYTHONHASHSEED"] = str(self.seed)

        # PyTorch 시드 고정
        try:
            import torch
            torch.manual_seed(self.seed)
            torch.cuda.manual_seed_all(self.seed)
            torch.backends.cudnn.deterministic = True
            torch.backends.cudnn.benchmark = False
            self.env_snapshot["torch_version"] = torch.__version__
            self.env_snapshot["cuda_available"] = (
                torch.cuda.is_available()
            )
        except ImportError:
            pass

        # TensorFlow 시드 고정
        try:
            import tensorflow as tf
            tf.random.set_seed(self.seed)
            self.env_snapshot["tf_version"] = tf.__version__
        except ImportError:
            pass

        self.experiment_log["seed"] = self.seed
        print(f"모든 랜덤 시드가 {self.seed}로 고정되었습니다.")

    def capture_environment(self) -> dict:
        """실행 환경 스냅샷을 캡처"""
        import sys
        self.env_snapshot.update({
            "python_version": sys.version,
            "platform": platform.platform(),
            "processor": platform.processor(),
            "timestamp": datetime.datetime.now().isoformat(),
        })
        # pip freeze로 패키지 목록 캡처
        try:
            result = subprocess.run(
                ["pip", "freeze"],
                capture_output=True, text=True, timeout=30
            )
            self.env_snapshot["packages"] = (
                result.stdout.strip().split("\n")
            )
        except Exception:
            self.env_snapshot["packages"] = []
        # Git 커밋 해시 기록
        try:
            result = subprocess.run(
                ["git", "rev-parse", "HEAD"],
                capture_output=True, text=True, timeout=10
            )
            self.env_snapshot["git_commit"] = result.stdout.strip()
        except Exception:
            self.env_snapshot["git_commit"] = "N/A"
        return self.env_snapshot

    def compute_data_hash(self, data_path: str) -> str:
        """데이터 파일의 해시값을 계산하여 버전 식별"""
        sha256 = hashlib.sha256()
        with open(data_path, "rb") as f:
            for chunk in iter(lambda: f.read(8192), b""):
                sha256.update(chunk)
        data_hash = sha256.hexdigest()[:12]
        self.experiment_log["data_hash"] = data_hash
        return data_hash

    def log_hyperparameters(self, params: Dict[str, Any]):
        """하이퍼파라미터를 기록"""
        self.experiment_log["hyperparameters"] = params

    def log_metrics(self, metrics: Dict[str, float]):
        """실험 결과 지표를 기록"""
        self.experiment_log["metrics"] = metrics

    def save_experiment_record(self, path: str):
        """실험 기록을 JSON으로 저장"""
        record = {
            "environment": self.env_snapshot,
            "experiment": self.experiment_log,
        }
        with open(path, "w", encoding="utf-8") as f:
            json.dump(record, f, ensure_ascii=False, indent=2)
        print(f"실험 기록 저장 완료: {path}")

# 사용 예시
manager = ReproducibilityManager(seed=42)
manager.fix_all_seeds()
manager.capture_environment()
manager.log_hyperparameters({
    "learning_rate": 0.001,
    "batch_size": 64,
    "epochs": 100,
    "model_type": "XGBoost",
})
manager.log_metrics({
    "accuracy": 0.8734,
    "f1_score": 0.8651,
    "auc": 0.9102,
})
manager.save_experiment_record("experiment_001.json")
\end{lstlisting}

%----------------------------------------------------------
\section{공정성 검증}
\label{sec:10.3}
%----------------------------------------------------------

AI 시스템이 특정 집단에 체계적으로 불리한 결과를 초래한다면, 이는 기술적 결함을 넘어 사회적, 법적 문제가 된다. 공정성(Fairness) 검증은 모델 개발 단계에서 반드시 수행해야 하는 핵심 거버넌스 활동이다.

\subsection{공정성 지표 체계}
\label{sec:10.3.1}

공정성을 측정하는 단일한 지표는 존재하지 않는다. 상황과 맥락에 따라 적절한 지표를 선택해야 하며, 주요 공정성 지표는 다음과 같다.

\paragraph{인구통계학적 동등성(Demographic Parity, DP)}
각 민감 그룹이 동일한 비율로 긍정적 결과(예: 대출 승인)를 받아야 한다.
\[
P(\hat{Y}=1 | A=a) = P(\hat{Y}=1 | A=b), \quad \forall a, b
\]
여기서 $A$는 민감 속성(예: 성별), $\hat{Y}$는 모델의 예측 결과이다.

\paragraph{균등 기회(Equalized Odds, EO)}
실제 긍정 사례에서 각 그룹의 진양성률(True Positive Rate)이 동일하고, 실제 부정 사례에서 각 그룹의 위양성률(False Positive Rate)이 동일해야 한다.
\[
P(\hat{Y}=1 | Y=y, A=a) = P(\hat{Y}=1 | Y=y, A=b), \quad \forall y, a, b
\]

\paragraph{보정(Calibration)}
각 그룹에서 모델이 예측한 확률이 실제 발생 확률과 일치해야 한다. 예를 들어, 모델이 70\% 확률로 부도를 예측한 경우, 남성 그룹과 여성 그룹 모두에서 실제로 약 70\%가 부도 나야 한다.

\paragraph{개인 공정성(Individual Fairness)}
유사한 특성을 가진 개인은 유사한 결과를 받아야 한다. 이는 ``비슷한 사람은 비슷하게 취급''이라는 직관에 기반한다.

\begin{table}[htbp]
\centering
\caption{공정성 지표별 적용 시나리오}
\label{tab:fairness_metrics}
\begin{tabularx}{\textwidth}{p{2.8cm}p{3cm}X}
\toprule
\textbf{지표} & \textbf{적합 시나리오} & \textbf{한계점} \\
\midrule
인구통계학적 동등성(DP) & 채용, 광고 노출 등 결과의 균등 배분이 중요한 경우 & 기저율(Base Rate) 차이를 무시. 정확도를 크게 희생할 수 있음 \\
\midrule
균등 기회(EO) & 대출 심사, 형사 재범 예측 등 정확한 분류가 중요한 경우 & 긍정 결과와 부정 결과의 오류율을 동시에 맞추기 어려움 \\
\midrule
보정(Calibration) & 의료 진단, 리스크 스코어링 등 확률 추정이 중요한 경우 & 그룹 간 기저율이 다르면 DP, EO와 동시 충족 불가 \\
\midrule
개인 공정성 & 개인화된 서비스, 맞춤형 추천 등 개별 판단이 중요한 경우 & ``유사성'' 정의가 주관적이며 측정이 어려움 \\
\bottomrule
\end{tabularx}
\end{table}

\subsection{공정성-정확도 트레이드오프 관리}
\label{sec:10.3.2}

공정성 제약을 강화하면 전체 정확도가 하락하는 트레이드오프가 발생한다. 이는 수학적으로 불가피한 경우가 많다---Chouldechova(2017)가 증명했듯이, 기저율이 다른 그룹에서 보정, 양성 예측도, 위양성률을 동시에 동등하게 만드는 것은 불가능하다.

거버넌스 관점에서 중요한 것은 이 트레이드오프를 투명하게 관리하는 것이다.

\begin{itemize}
  \item \textbf{트레이드오프 시각화}: 공정성 지표와 정확도 지표의 파레토 프런티어(Pareto Frontier)를 시각화하여 의사결정자에게 제시한다.
  \item \textbf{허용 임계값 사전 결정}: 비즈니스 오너, 법무, 윤리 담당자가 합의하여 공정성 지표의 최소 기준과 정확도의 최소 기준을 사전에 정의한다.
  \item \textbf{의사결정 기록}: 어떤 트레이드오프 지점을 선택했는지, 그 이유는 무엇인지를 문서화한다.
\end{itemize}

\subsection{편향 완화 기법}
\label{sec:10.3.3}

편향(Bias)은 모델 개발 파이프라인의 여러 단계에서 개입할 수 있으며, 각 단계에 맞는 완화 기법이 존재한다.

\paragraph{전처리 기법(Pre-processing)}
학습 데이터 자체의 편향을 제거하는 접근이다.
\begin{itemize}
  \item \textbf{리웨이팅(Reweighting)}: 소수 그룹의 샘플에 더 높은 가중치를 부여
  \item \textbf{리샘플링(Resampling)}: 과소 대표 그룹을 오버샘플링하거나 과대 대표 그룹을 언더샘플링
  \item \textbf{디스패럿 임팩트 리무버(Disparate Impact Remover)}: 민감 속성의 영향을 데이터에서 제거하면서 순위를 보존
\end{itemize}

\paragraph{중간처리 기법(In-processing)}
학습 알고리즘 자체에 공정성 제약을 포함하는 접근이다.
\begin{itemize}
  \item \textbf{적대적 디바이어싱(Adversarial Debiasing)}: 적대적 네트워크를 사용하여 모델이 민감 속성을 추론하지 못하도록 학습
  \item \textbf{제약 최적화(Constrained Optimization)}: 공정성 지표를 제약 조건으로 추가한 최적화 문제로 변환
  \item \textbf{정규화(Regularization)}: 손실 함수에 공정성 패널티 항을 추가
\end{itemize}

\paragraph{후처리 기법(Post-processing)}
학습된 모델의 출력을 조정하는 접근이다.
\begin{itemize}
  \item \textbf{임계값 조정(Threshold Adjustment)}: 각 그룹별로 최적의 분류 임계값을 다르게 설정
  \item \textbf{보정된 균등 기회(Calibrated Equalized Odds)}: 예측 확률을 보정하여 균등 기회를 달성
  \item \textbf{거부 옵션 분류(Reject Option Classification)}: 결정 경계 근처의 불확실한 예측에 대해 소수 그룹에 유리하게 판정
\end{itemize}

\subsection{공정성 자동 검증 파이프라인}
\label{sec:10.3.4}

다음 코드는 Fairlearn 라이브러리를 기반으로 공정성 지표를 자동으로 계산하고 보고하는 파이프라인을 보여준다.

\begin{lstlisting}[language=Python, caption={Fairlearn 기반 공정성 자동 검증 파이프라인}, label={lst:fairness_pipeline}]
import numpy as np
from dataclasses import dataclass, field
from typing import Dict, List, Tuple, Optional

@dataclass
class FairnessReport:
    """공정성 검증 보고서"""
    model_name: str
    sensitive_feature_name: str
    metrics: Dict[str, Dict] = field(default_factory=dict)
    overall_pass: bool = False
    violations: List[str] = field(default_factory=list)

class FairnessValidationPipeline:
    """공정성 자동 검증 파이프라인"""

    def __init__(
        self,
        model_name: str,
        dp_threshold: float = 0.8,
        eo_threshold: float = 0.1,
    ):
        self.model_name = model_name
        self.dp_threshold = dp_threshold  # 4/5 규칙
        self.eo_threshold = eo_threshold
        self.reports: List[FairnessReport] = []

    def compute_demographic_parity(
        self,
        y_pred: np.ndarray,
        sensitive: np.ndarray,
    ) -> Dict:
        """인구통계학적 동등성 계산"""
        groups = np.unique(sensitive)
        selection_rates = {}
        for g in groups:
            mask = sensitive == g
            selection_rates[str(g)] = float(
                np.mean(y_pred[mask])
            )
        rates = list(selection_rates.values())
        max_rate = max(rates)
        min_rate = min(rates)
        ratio = min_rate / max_rate if max_rate > 0 else 0
        return {
            "selection_rates": selection_rates,
            "ratio": round(ratio, 4),
            "pass": ratio >= self.dp_threshold,
        }

    def compute_equalized_odds(
        self,
        y_true: np.ndarray,
        y_pred: np.ndarray,
        sensitive: np.ndarray,
    ) -> Dict:
        """균등 기회 계산"""
        groups = np.unique(sensitive)
        tpr_by_group = {}
        fpr_by_group = {}
        for g in groups:
            mask = sensitive == g
            pos_mask = mask & (y_true == 1)
            neg_mask = mask & (y_true == 0)
            tpr = (
                np.mean(y_pred[pos_mask])
                if pos_mask.sum() > 0 else 0
            )
            fpr = (
                np.mean(y_pred[neg_mask])
                if neg_mask.sum() > 0 else 0
            )
            tpr_by_group[str(g)] = round(float(tpr), 4)
            fpr_by_group[str(g)] = round(float(fpr), 4)
        tpr_diff = max(tpr_by_group.values()) - min(
            tpr_by_group.values()
        )
        fpr_diff = max(fpr_by_group.values()) - min(
            fpr_by_group.values()
        )
        return {
            "tpr_by_group": tpr_by_group,
            "fpr_by_group": fpr_by_group,
            "tpr_difference": round(tpr_diff, 4),
            "fpr_difference": round(fpr_diff, 4),
            "pass": (
                tpr_diff <= self.eo_threshold
                and fpr_diff <= self.eo_threshold
            ),
        }

    def validate(
        self,
        y_true: np.ndarray,
        y_pred: np.ndarray,
        sensitive_features: Dict[str, np.ndarray],
    ) -> List[FairnessReport]:
        """전체 공정성 검증 실행"""
        self.reports = []
        for feat_name, feat_values in sensitive_features.items():
            report = FairnessReport(
                model_name=self.model_name,
                sensitive_feature_name=feat_name,
            )
            dp = self.compute_demographic_parity(
                y_pred, feat_values
            )
            report.metrics["demographic_parity"] = dp
            if not dp["pass"]:
                report.violations.append(
                    f"DP 위반: 비율 {dp['ratio']} "
                    f"< 임계값 {self.dp_threshold}"
                )
            eo = self.compute_equalized_odds(
                y_true, y_pred, feat_values
            )
            report.metrics["equalized_odds"] = eo
            if not eo["pass"]:
                report.violations.append(
                    f"EO 위반: TPR 차이 {eo['tpr_difference']}"
                    f", FPR 차이 {eo['fpr_difference']}"
                )
            report.overall_pass = len(report.violations) == 0
            self.reports.append(report)

        return self.reports

    def print_summary(self):
        """검증 결과 요약 출력"""
        for report in self.reports:
            status = "통과" if report.overall_pass else "실패"
            print(f"\n=== {report.sensitive_feature_name} ===")
            print(f"  상태: {status}")
            for name, result in report.metrics.items():
                pass_str = "통과" if result["pass"] else "실패"
                print(f"  {name}: {pass_str}")
            for v in report.violations:
                print(f"  [위반] {v}")

# 사용 예시
pipeline = FairnessValidationPipeline(
    model_name="CreditScoring-v2",
    dp_threshold=0.8,
    eo_threshold=0.1,
)
# 모의 데이터
np.random.seed(42)
y_true = np.random.randint(0, 2, 1000)
y_pred = np.random.randint(0, 2, 1000)
gender = np.random.choice(["M", "F"], 1000)
age_group = np.random.choice(["young", "middle", "senior"], 1000)

reports = pipeline.validate(
    y_true, y_pred,
    {"gender": gender, "age_group": age_group},
)
pipeline.print_summary()
\end{lstlisting}

%----------------------------------------------------------
\section{강건성과 보안 테스트}
\label{sec:10.4}
%----------------------------------------------------------

모델이 정상적인 입력에서만 잘 작동하는 것은 충분하지 않다. 악의적인 공격이나 예상치 못한 입력에 대해서도 안정적으로 동작하는 강건성(Robustness)은 프로덕션 배포의 필수 조건이다\cite{goodfellow2014explaining}.

\subsection{적대적 공격 유형과 방어 기법}
\label{sec:10.4.1}

적대적 공격(Adversarial Attack)은 모델의 입력을 미세하게 조작하여 오분류를 유도하는 기법이다. \cite{goodfellow2014explaining}가 제안한 FGSM(Fast Gradient Sign Method)은 이 분야의 시초가 되었다.

\begin{table}[htbp]
\centering
\caption{주요 적대적 공격 유형과 방어 기법}
\label{tab:adversarial_attacks}
\begin{tabularx}{\textwidth}{p{2.5cm}p{3.5cm}X}
\toprule
\textbf{공격 유형} & \textbf{설명} & \textbf{방어 기법} \\
\midrule
FGSM & 그래디언트 방향으로 단일 스텝 섭동 추가. 빠르지만 간단한 공격 & 적대적 학습(Adversarial Training)으로 강건성 향상 \\
\midrule
PGD & 반복적 그래디언트 기반 공격. FGSM보다 강력한 공격 & 인증된 방어(Certified Defense), 무작위 평활화(Randomized Smoothing) \\
\midrule
C\&W & 최적화 기반 공격. 최소 섭동으로 최대 효과 & 방어 증류(Defensive Distillation), 입력 변환(Input Transformation) \\
\midrule
데이터 중독(Data Poisoning) & 학습 데이터에 악의적 샘플을 주입 & 데이터 검증, 이상치 탐지, 강건 통계 기법 \\
\midrule
백도어 공격(Backdoor Attack) & 특정 트리거 패턴 입력 시 오동작 유도 & 뉴런 가지치기(Neural Pruning), 미세조정(Fine-tuning) \\
\midrule
모델 역전(Model Inversion) & 모델 출력에서 학습 데이터 복원 시도 & 차분 프라이버시(Differential Privacy), 출력 제한 \\
\bottomrule
\end{tabularx}
\end{table}

\subsection{입력 검증과 이상치 탐지}
\label{sec:10.4.2}

프로덕션 환경에서 모델은 학습 시 보지 못한 다양한 입력을 받게 된다. 분포 외(Out-of-Distribution, OOD) 입력을 탐지하고 적절히 처리하는 메커니즘이 필요하다.

\begin{itemize}
  \item \textbf{입력 범위 검증}: 각 특성(Feature)의 유효 범위를 정의하고, 범위를 벗어난 입력을 감지
  \item \textbf{통계적 이상치 탐지}: 학습 데이터의 분포와 비교하여 이상 입력을 식별 (예: 마할라노비스 거리(Mahalanobis Distance), 격리 숲(Isolation Forest))
  \item \textbf{신뢰도 기반 거부}: 모델의 예측 신뢰도가 임계값 이하인 경우 자동 판단을 보류하고 사람의 검토를 요청
  \item \textbf{입력 무결성 검증}: 입력 데이터의 형식, 타입, 필수 필드 존재 여부를 검증
\end{itemize}

\subsection{모델 추출 공격 방어}
\label{sec:10.4.3}

모델 추출 공격(Model Extraction Attack)은 공격자가 API를 통해 대량의 쿼리를 보내 모델을 복제하려는 시도이다. 이는 지적 재산 보호와 보안 모두에 위협이 된다.

\begin{itemize}
  \item \textbf{쿼리 제한(Rate Limiting)}: API 호출 빈도와 총량을 제한
  \item \textbf{워터마킹(Watermarking)}: 모델 출력에 식별 가능한 패턴을 삽입하여 복제 모델을 추적
  \item \textbf{출력 정밀도 제한}: 예측 확률의 소수점 자릿수를 제한하여 정보 노출을 최소화
  \item \textbf{쿼리 패턴 모니터링}: 비정상적인 쿼리 패턴(예: 체계적인 입력 공간 탐색)을 탐지하여 차단
\end{itemize}

\subsection{적대적 강건성 테스트 프레임워크}
\label{sec:10.4.4}

다음 코드는 모델의 적대적 강건성을 체계적으로 테스트하는 프레임워크를 보여준다.

\begin{lstlisting}[language=Python, caption={적대적 강건성 테스트 프레임워크}, label={lst:adversarial_test}]
import numpy as np
from dataclasses import dataclass, field
from typing import Dict, List, Callable, Any, Optional

@dataclass
class RobustnessTestResult:
    """강건성 테스트 결과"""
    test_name: str
    original_accuracy: float
    perturbed_accuracy: float
    robustness_score: float  # 0~1
    passed: bool
    details: Dict[str, Any] = field(default_factory=dict)

class AdversarialRobustnessFramework:
    """적대적 강건성 테스트 프레임워크"""

    def __init__(
        self,
        model,
        robustness_threshold: float = 0.8,
    ):
        self.model = model
        self.threshold = robustness_threshold
        self.results: List[RobustnessTestResult] = []

    def _predict(self, X: np.ndarray) -> np.ndarray:
        """모델 예측 래퍼"""
        return self.model.predict(X)

    def test_gaussian_noise(
        self,
        X: np.ndarray,
        y: np.ndarray,
        noise_levels: List[float] = None,
    ) -> RobustnessTestResult:
        """가우시안 노이즈에 대한 강건성 테스트"""
        if noise_levels is None:
            noise_levels = [0.01, 0.05, 0.1, 0.2]
        original_pred = self._predict(X)
        orig_acc = float(np.mean(original_pred == y))
        noise_results = {}
        for sigma in noise_levels:
            noise = np.random.normal(0, sigma, X.shape)
            X_noisy = X + noise
            noisy_pred = self._predict(X_noisy)
            noisy_acc = float(np.mean(noisy_pred == y))
            noise_results[f"sigma_{sigma}"] = round(
                noisy_acc, 4
            )
        worst_acc = min(noise_results.values())
        robustness = worst_acc / orig_acc if orig_acc > 0 else 0
        result = RobustnessTestResult(
            test_name="가우시안 노이즈 테스트",
            original_accuracy=round(orig_acc, 4),
            perturbed_accuracy=round(worst_acc, 4),
            robustness_score=round(robustness, 4),
            passed=robustness >= self.threshold,
            details=noise_results,
        )
        self.results.append(result)
        return result

    def test_feature_perturbation(
        self,
        X: np.ndarray,
        y: np.ndarray,
        feature_idx: int,
        perturbation_pct: float = 0.1,
    ) -> RobustnessTestResult:
        """특정 특성 섭동에 대한 강건성 테스트"""
        original_pred = self._predict(X)
        orig_acc = float(np.mean(original_pred == y))
        X_perturbed = X.copy()
        feature_std = np.std(X[:, feature_idx])
        perturbation = feature_std * perturbation_pct
        X_perturbed[:, feature_idx] += perturbation
        perturbed_pred = self._predict(X_perturbed)
        pert_acc = float(np.mean(perturbed_pred == y))
        robustness = pert_acc / orig_acc if orig_acc > 0 else 0
        result = RobustnessTestResult(
            test_name=f"특성[{feature_idx}] 섭동 테스트",
            original_accuracy=round(orig_acc, 4),
            perturbed_accuracy=round(pert_acc, 4),
            robustness_score=round(robustness, 4),
            passed=robustness >= self.threshold,
            details={
                "feature_index": feature_idx,
                "perturbation_pct": perturbation_pct,
            },
        )
        self.results.append(result)
        return result

    def test_input_validation(
        self,
        X: np.ndarray,
        valid_ranges: Dict[int, tuple],
    ) -> RobustnessTestResult:
        """입력 범위 검증 테스트"""
        violations = {}
        total_checks = 0
        total_violations = 0
        for feat_idx, (low, high) in valid_ranges.items():
            col = X[:, feat_idx]
            out_of_range = np.sum(
                (col < low) | (col > high)
            )
            total_checks += len(col)
            total_violations += out_of_range
            if out_of_range > 0:
                violations[f"feature_{feat_idx}"] = int(
                    out_of_range
                )
        violation_rate = (
            total_violations / total_checks
            if total_checks > 0 else 0
        )
        result = RobustnessTestResult(
            test_name="입력 범위 검증",
            original_accuracy=1.0,
            perturbed_accuracy=round(1 - violation_rate, 4),
            robustness_score=round(1 - violation_rate, 4),
            passed=violation_rate < 0.01,  # 1% 미만
            details={"violations": violations},
        )
        self.results.append(result)
        return result

    def generate_report(self) -> Dict:
        """전체 강건성 테스트 보고서 생성"""
        all_passed = all(r.passed for r in self.results)
        report = {
            "overall_status": "통과" if all_passed else "실패",
            "total_tests": len(self.results),
            "passed_tests": sum(
                1 for r in self.results if r.passed
            ),
            "details": [],
        }
        for r in self.results:
            report["details"].append({
                "test": r.test_name,
                "status": "통과" if r.passed else "실패",
                "robustness_score": r.robustness_score,
                "original_acc": r.original_accuracy,
                "perturbed_acc": r.perturbed_accuracy,
            })
        return report
\end{lstlisting}

%----------------------------------------------------------
\section{설명가능성(XAI) 구현}
\label{sec:10.5}
%----------------------------------------------------------

``AI가 왜 대출을 거절했습니까?''라는 고객의 질문에 ``알고리즘이 결정했습니다''라고 답할 수는 없다. 특히 금융, 공공, 의료 등 고위험 분야에서는 설명 요구권(Right to Explanation)이 법적으로 보장되는 추세다. \cite{arrieta2020explainable}는 책임 있는 AI를 위해 설명가능한 AI(eXplainable AI, XAI)가 필수적임을 강조한다.

\subsection{글로벌 vs 로컬 설명가능성}
\label{sec:10.5.1}

설명가능성은 그 범위에 따라 글로벌(Global)과 로컬(Local) 수준으로 나뉜다.

\paragraph{글로벌 설명가능성}
모델 전체의 행동 패턴을 설명한다. ``이 모델은 전반적으로 어떤 특성을 중요하게 여기는가?''에 답한다.
\begin{itemize}
  \item \textbf{특성 중요도(Feature Importance)}: 모델의 전반적 결정에 각 특성이 기여하는 정도
  \item \textbf{SHAP\cite{lundberg2017unified} 요약 플롯(SHAP Summary Plot)}: 모든 샘플에 걸쳐 각 특성의 SHAP 값 분포를 시각화
  \item \textbf{부분 의존도 플롯(Partial Dependence Plot, PDP)}: 특정 특성의 값 변화에 따른 모델 출력의 평균적 변화
  \item \textbf{글로벌 대리 모델(Global Surrogate)}: 복잡한 모델을 해석 가능한 단순 모델(예: 결정 트리)로 근사
\end{itemize}

\paragraph{로컬 설명가능성}
개별 예측에 대한 설명을 제공한다. ``이 특정 고객의 대출은 왜 거절되었는가?''에 답한다.
\begin{itemize}
  \item \textbf{LIME\cite{ribeiro2016why}(Local Interpretable Model-agnostic Explanations)}: 개별 예측 주변에서 해석 가능한 모델을 학습하여 설명 생성
  \item \textbf{개별 SHAP 값(Individual SHAP Values)}: 특정 예측에 대한 각 특성의 기여도를 분해
  \item \textbf{반사실적 설명(Counterfactual Explanation)}: ``어떤 입력이 달라졌다면 결과가 바뀌었을 것''이라는 형태의 설명
  \item \textbf{Grad-CAM(Gradient-weighted Class Activation Mapping)}: 이미지 분류에서 모델이 주목한 영역을 시각화
\end{itemize}

\subsection{주요 설명 기법 비교}
\label{sec:10.5.2}

\begin{table}[htbp]
\centering
\caption{SHAP, LIME, Grad-CAM 비교}
\label{tab:xai_comparison}
\begin{tabularx}{\textwidth}{p{2cm}p{2.3cm}p{2.3cm}p{2.3cm}X}
\toprule
\textbf{기준} & \textbf{SHAP} & \textbf{LIME} & \textbf{Grad-CAM} & \textbf{반사실적 설명} \\
\midrule
범위 & 글로벌 + 로컬 & 로컬 & 로컬 & 로컬 \\
\midrule
모델 의존성 & 모델 비의존적 (KernelSHAP) 또는 모델 특화 (TreeSHAP) & 모델 비의존적 & CNN 특화 & 모델 비의존적 \\
\midrule
데이터 유형 & 정형, 텍스트, 이미지 & 정형, 텍스트, 이미지 & 이미지 & 정형 \\
\midrule
계산 비용 & 중간\textasciitilde{}높음 & 중간 & 낮음 & 높음 \\
\midrule
이론적 기반 & 샤플리 값(게임 이론) & 지역 선형 근사 & 그래디언트 기반 활성화 & 인과 추론 \\
\midrule
장점 & 수학적 보장, 일관된 특성 기여도 & 직관적, 빠른 구현 & 시각적으로 명확 & 사용자 친화적 설명 \\
\midrule
단점 & 높은 계산 비용 (KernelSHAP) & 불안정성(같은 입력에 다른 설명 가능) & CNN에만 적용 & 유효한 반사실 생성이 어려움 \\
\bottomrule
\end{tabularx}
\end{table}

\subsection{반사실적 설명 구현}
\label{sec:10.5.3}

반사실적 설명(Counterfactual Explanation)은 ``연 소득이 10\%만 더 높았다면 대출이 승인되었을 것입니다''와 같은 형태로, 사용자에게 가장 직관적이고 실행 가능한 피드백을 제공한다.

효과적인 반사실적 설명은 다음 조건을 만족해야 한다.
\begin{itemize}
  \item \textbf{유효성(Validity)}: 반사실적 입력이 실제로 다른 결과를 이끌어내야 한다.
  \item \textbf{근접성(Proximity)}: 원본 입력과 최소한의 차이만 가져야 한다.
  \item \textbf{실현가능성(Plausibility)}: 제안된 변경이 현실적으로 가능해야 한다 (예: ``나이를 10살 줄이세요''는 비현실적).
  \item \textbf{희소성(Sparsity)}: 가능한 적은 수의 특성만 변경해야 한다.
\end{itemize}

\subsection{설명가능성 수준 결정 프레임워크}
\label{sec:10.5.4}

모든 모델에 동일한 수준의 설명가능성을 요구하는 것은 비효율적이다. 리스크 등급에 따라 요구 수준을 차등 적용한다.

\begin{table}[htbp]
\centering
\caption{리스크 등급별 설명가능성 요구 수준}
\label{tab:xai_levels}
\begin{tabularx}{\textwidth}{p{2cm}p{3cm}X}
\toprule
\textbf{리스크 등급} & \textbf{요구 수준} & \textbf{상세 요구사항} \\
\midrule
고위험 & 완전한 설명 & 글로벌 + 로컬 설명 필수. 반사실적 설명 제공. 설명의 정확도 검증. 비전문가도 이해 가능한 자연어 설명. 규제 감사 대응 문서화 \\
\midrule
중위험 & 표준 설명 & 글로벌 특성 중요도 + 요청 시 로컬 설명 제공. SHAP 또는 LIME 기반 기여도 분석. 내부 감사 대응 수준 문서화 \\
\midrule
저위험 & 기본 설명 & 글로벌 특성 중요도 제공. 이상 탐지 시 로컬 설명 생성. 기술 문서 수준 기록 \\
\midrule
최저위험 & 모니터링 & 모델 성능 대시보드. 문제 발생 시 사후 분석 가능한 로깅 체계 \\
\bottomrule
\end{tabularx}
\end{table}

\subsection{통합 설명가능성 보고서 생성기}
\label{sec:10.5.5}

다음 코드는 여러 XAI 기법을 통합하여 하나의 설명가능성 보고서를 자동 생성하는 시스템을 보여준다.

\begin{lstlisting}[language=Python, caption={통합 설명가능성 보고서 생성기}, label={lst:xai_report}]
import json
import numpy as np
from dataclasses import dataclass, field
from typing import Dict, List, Any, Optional

@dataclass
class ExplanationResult:
    """단일 설명 결과"""
    method: str
    scope: str  # "global" or "local"
    feature_contributions: Dict[str, float]
    description: str = ""

@dataclass
class XAIReport:
    """설명가능성 보고서"""
    model_name: str
    risk_level: str
    explanations: List[ExplanationResult] = field(
        default_factory=list
    )
    counterfactuals: List[Dict] = field(default_factory=list)
    metadata: Dict[str, Any] = field(default_factory=dict)

class IntegratedXAIReportGenerator:
    """통합 설명가능성 보고서 생성기"""

    def __init__(
        self,
        model,
        model_name: str,
        risk_level: str = "high",
        feature_names: List[str] = None,
    ):
        self.model = model
        self.model_name = model_name
        self.risk_level = risk_level
        self.feature_names = feature_names or []
        self.report = XAIReport(
            model_name=model_name,
            risk_level=risk_level,
        )

    def compute_permutation_importance(
        self,
        X: np.ndarray,
        y: np.ndarray,
        n_repeats: int = 10,
    ) -> ExplanationResult:
        """순열 중요도 기반 글로벌 설명"""
        from sklearn.inspection import permutation_importance
        result = permutation_importance(
            self.model, X, y,
            n_repeats=n_repeats, random_state=42,
        )
        contributions = {}
        for i, imp in enumerate(result.importances_mean):
            name = (
                self.feature_names[i]
                if i < len(self.feature_names)
                else f"feature_{i}"
            )
            contributions[name] = round(float(imp), 4)
        # 중요도 내림차순 정렬
        contributions = dict(
            sorted(
                contributions.items(),
                key=lambda x: abs(x[1]),
                reverse=True,
            )
        )
        explanation = ExplanationResult(
            method="Permutation Importance",
            scope="global",
            feature_contributions=contributions,
            description="순열 중요도 기반 글로벌 특성 기여도",
        )
        self.report.explanations.append(explanation)
        return explanation

    def compute_shap_local(
        self,
        X_instance: np.ndarray,
        X_background: np.ndarray,
    ) -> ExplanationResult:
        """SHAP 기반 로컬 설명"""
        try:
            import shap
            explainer = shap.KernelExplainer(
                self.model.predict_proba, X_background[:100]
            )
            shap_values = explainer.shap_values(
                X_instance.reshape(1, -1)
            )
            if isinstance(shap_values, list):
                values = shap_values[1][0]
            else:
                values = shap_values[0]
        except ImportError:
            # SHAP 미설치 시 근사 계산
            values = np.random.randn(len(X_instance)) * 0.1
        contributions = {}
        for i, val in enumerate(values):
            name = (
                self.feature_names[i]
                if i < len(self.feature_names)
                else f"feature_{i}"
            )
            contributions[name] = round(float(val), 4)
        explanation = ExplanationResult(
            method="SHAP",
            scope="local",
            feature_contributions=contributions,
            description="SHAP 기반 개별 예측 기여도 분해",
        )
        self.report.explanations.append(explanation)
        return explanation

    def generate_counterfactual(
        self,
        X_instance: np.ndarray,
        desired_outcome: int,
        mutable_features: List[int] = None,
        step_size: float = 0.1,
        max_iter: int = 100,
    ) -> Optional[Dict]:
        """간단한 반사실적 설명 생성"""
        if mutable_features is None:
            mutable_features = list(range(len(X_instance)))
        x_cf = X_instance.copy()
        for iteration in range(max_iter):
            pred = self.model.predict(x_cf.reshape(1, -1))[0]
            if pred == desired_outcome:
                changes = {}
                for i in mutable_features:
                    diff = x_cf[i] - X_instance[i]
                    if abs(diff) > 1e-6:
                        name = (
                            self.feature_names[i]
                            if i < len(self.feature_names)
                            else f"feature_{i}"
                        )
                        changes[name] = {
                            "original": round(
                                float(X_instance[i]), 4
                            ),
                            "counterfactual": round(
                                float(x_cf[i]), 4
                            ),
                            "change": round(float(diff), 4),
                        }
                cf_result = {
                    "found": True,
                    "iterations": iteration + 1,
                    "changes": changes,
                }
                self.report.counterfactuals.append(cf_result)
                return cf_result
            # 랜덤 방향으로 탐색
            for feat in mutable_features:
                x_cf[feat] += np.random.choice(
                    [-step_size, step_size]
                ) * np.std(X_instance)
        return {"found": False, "iterations": max_iter}

    def generate_full_report(
        self, output_path: str = None
    ) -> Dict:
        """전체 보고서를 생성하고 저장"""
        report_dict = {
            "model_name": self.report.model_name,
            "risk_level": self.report.risk_level,
            "explanations": [
                {
                    "method": e.method,
                    "scope": e.scope,
                    "description": e.description,
                    "top_features": dict(
                        list(e.feature_contributions.items())[:5]
                    ),
                }
                for e in self.report.explanations
            ],
            "counterfactuals": self.report.counterfactuals,
        }
        if output_path:
            with open(output_path, "w", encoding="utf-8") as f:
                json.dump(
                    report_dict, f, ensure_ascii=False, indent=2
                )
        return report_dict
\end{lstlisting}

%----------------------------------------------------------
\section{모델 승인과 릴리스 관리}
\label{sec:10.6}
%----------------------------------------------------------

개발과 검증이 끝난 모델이 실제 운영 환경(프로덕션(Production))으로 나가기 전, 마지막 관문(Gate)을 거쳐야 한다. 이 관문은 기술적 품질뿐만 아니라 윤리적, 법적, 보안적 측면을 종합적으로 검증하는 과정이다.

\subsection{모델 승인 게이트(다단계 리뷰)}
\label{sec:10.6.1}

고위험 모델의 경우, 단일 승인이 아닌 다단계 리뷰 프로세스를 거쳐야 한다.

\paragraph{1단계: 기술 리뷰(Technical Review)}
데이터 사이언스 팀 내부에서 수행하는 기술적 검증이다.
\begin{itemize}
  \item 모델 성능이 사전 정의된 KPI를 달성하는지 확인
  \item 코드 리뷰 완료 여부 확인
  \item 재현성 테스트 통과 여부 확인
  \item 단위 테스트 및 통합 테스트 통과 여부 확인
\end{itemize}

\paragraph{2단계: 공정성 및 윤리 리뷰(Fairness \& Ethics Review)}
AI 윤리 담당자 또는 공정성 검증팀이 수행한다.
\begin{itemize}
  \item 공정성 지표가 사전 정의된 임계값을 충족하는지 확인
  \item 편향 완화 조치의 적절성 평가
  \item 설명가능성 수준이 리스크 등급에 맞는지 확인
  \item 윤리적 영향 평가서(Ethical Impact Assessment) 검토
\end{itemize}

\paragraph{3단계: 법무 및 컴플라이언스 리뷰(Legal \& Compliance Review)}
법무팀과 규제 대응팀이 수행한다.
\begin{itemize}
  \item 관련 법규(개인정보보호법, AI 기본법 등) 준수 여부 확인
  \item 데이터 사용 동의 범위 적합성 확인
  \item 설명 요구권 대응 가능 여부 확인
  \item 필요한 고지 사항 확인
\end{itemize}

\paragraph{4단계: 보안 리뷰(Security Review)}
정보보안팀이 수행한다.
\begin{itemize}
  \item 적대적 공격에 대한 강건성 테스트 결과 확인
  \item 모델 파일 무결성 검증 체계 확인
  \item API 보안 설정(인증, 권한, 암호화) 확인
  \item 의존성 라이브러리 보안 취약점 스캔 결과 확인
\end{itemize}

\paragraph{5단계: 최종 승인(Final Approval)}
배포 승인 위원회(Model Review Board)에서 각 단계의 리뷰 결과를 종합하여 최종 결정을 내린다.
\begin{itemize}
  \item \textbf{승인(Go)}: 배포 승인.
  \item \textbf{반려(No-Go)}: 문제 발견으로 반려. 재학습 또는 수정 후 재상정.
  \item \textbf{조건부 승인(Conditional Go)}: 배포하되, 초기 트래픽을 10\%로 제한(카나리(Canary) 배포)하고 모니터링을 강화하는 조건.
\end{itemize}

\subsection{배포 전 체크리스트}
\label{sec:10.6.2}

\begin{table}[htbp]
\centering
\caption{배포 전 종합 체크리스트}
\label{tab:deployment_checklist}
\begin{tabularx}{\textwidth}{p{2cm}p{4cm}X}
\toprule
\textbf{영역} & \textbf{체크 항목} & \textbf{검증 기준} \\
\midrule
\multirow{3}{*}{기술} & 성능 KPI 달성 & 사전 합의된 모든 성능 지표가 임계값 이상 \\
\cmidrule{2-3}
& 재현성 확인 & 동일 데이터/설정으로 결과 재현 가능 \\
\cmidrule{2-3}
& 추론 성능 & 지연시간(Latency), 처리량(Throughput) 요구사항 충족 \\
\midrule
\multirow{2}{*}{윤리} & 공정성 검증 통과 & 모든 민감 속성에 대해 공정성 임계값 충족 \\
\cmidrule{2-3}
& 설명가능성 구현 & 리스크 등급에 맞는 설명 기능 구현 완료 \\
\midrule
\multirow{2}{*}{법무} & 규제 준수 확인 & 관련 법규 요구사항 충족 확인서 \\
\cmidrule{2-3}
& 데이터 동의 확인 & 학습 데이터의 사용 동의 범위 적합 \\
\midrule
\multirow{2}{*}{보안} & 강건성 테스트 통과 & 적대적 공격 테스트 통과 \\
\cmidrule{2-3}
& 취약점 스캔 통과 & 의존성 보안 취약점 없음 확인 \\
\midrule
운영 & 모니터링 설정 완료 & 성능, 공정성, 드리프트 모니터링 대시보드 구축 \\
\cmidrule{2-3}
& 롤백 절차 수립 & 비상 시 롤백 절차 문서화 및 테스트 완료 \\
\bottomrule
\end{tabularx}
\end{table}

\subsection{카나리 배포와 A/B 테스트 거버넌스}
\label{sec:10.6.3}

모델 배포는 빅뱅(Big-bang) 방식이 아닌 점진적 방식을 권장한다.

\paragraph{카나리 배포(Canary Deployment)}
새 모델을 전체 트래픽의 일부(예: 5\textasciitilde{}10\%)에만 먼저 적용하고, 문제가 없음을 확인한 후 점진적으로 확대하는 방식이다.
\begin{itemize}
  \item 카나리 비율과 확대 일정을 사전에 정의
  \item 카나리 기간 동안의 핵심 모니터링 지표(성능, 오류율, 지연시간) 정의
  \item 자동 롤백 트리거 조건 설정 (예: 오류율이 기존 대비 2배 이상 증가 시)
  \item 카나리 결과에 대한 승인 절차 (카나리 통과 시에만 전체 배포 진행)
\end{itemize}

\paragraph{A/B 테스트 거버넌스}
기존 모델과 신규 모델을 동시에 운영하며 성능을 비교하는 A/B 테스트는 과학적인 모델 교체 의사결정의 기반이 된다.
\begin{itemize}
  \item 테스트 설계: 표본 크기, 테스트 기간, 통계적 유의수준을 사전 정의
  \item 지표 정의: 주요 지표(Primary Metric)와 보조 지표(Guardrail Metric)를 구분하여 설정
  \item 윤리적 고려: A/B 테스트 대상이 특정 그룹에 불이익을 주지 않는지 확인
  \item 결과 해석: 통계적 유의성에 기반한 객관적 의사결정. p-해킹(p-hacking)이나 조기 중단(early stopping) 방지
\end{itemize}

\subsection{롤백 정책과 비상 절차}
\label{sec:10.6.4}

배포 후 심각한 문제가 발견될 경우, 신속하게 이전 버전으로 되돌릴 수 있어야 한다.

\paragraph{롤백 트리거}
다음 조건 중 하나라도 충족되면 즉시 롤백을 개시한다.
\begin{itemize}
  \item 핵심 성능 지표가 사전 정의된 임계값 아래로 하락
  \item 공정성 지표 위반이 탐지됨
  \item 보안 취약점이 발견됨
  \item 예상치 못한 시스템 오류 발생률이 급증
  \item 규제 기관의 시정 요구
\end{itemize}

\paragraph{롤백 절차}
\begin{enumerate}
  \item 이상 탐지 시 자동 알림 발송 (운영팀, 개발팀, 관리자)
  \item 1차 대응: 트래픽을 이전 모델로 전환 (자동 또는 수동)
  \item 영향 범위 분석: 롤백 전 기간 동안 영향받은 건수 파악
  \item 원인 분석: 문제의 근본 원인을 식별
  \item 사후 보고서(Post-mortem) 작성
  \item 재배포 계획 수립 및 승인
\end{enumerate}

\subsection{배포 승인 자동화 시스템}
\label{sec:10.6.5}

다음 코드는 다단계 승인 프로세스를 자동화하여 모든 게이트를 통과한 모델만 배포할 수 있도록 하는 시스템을 보여준다.

\begin{lstlisting}[language=Python, caption={배포 승인 자동화 시스템}, label={lst:deployment_approval}]
import json
import datetime
from dataclasses import dataclass, field
from typing import Dict, List, Optional
from enum import Enum

class ApprovalStatus(Enum):
    PENDING = "pending"
    APPROVED = "approved"
    REJECTED = "rejected"
    CONDITIONAL = "conditional"

@dataclass
class GateResult:
    """단일 승인 게이트 결과"""
    gate_name: str
    status: ApprovalStatus
    reviewer: str
    timestamp: str
    checks: Dict[str, bool] = field(default_factory=dict)
    comments: str = ""
    conditions: List[str] = field(default_factory=list)

class DeploymentApprovalSystem:
    """배포 승인 자동화 시스템"""

    def __init__(self, model_name: str, model_version: str):
        self.model_name = model_name
        self.model_version = model_version
        self.gates: List[GateResult] = []
        self.final_status = ApprovalStatus.PENDING

    def run_technical_gate(
        self,
        performance_metrics: Dict[str, float],
        performance_thresholds: Dict[str, float],
        reproducibility_passed: bool,
        latency_ms: float,
        latency_threshold: float,
    ) -> GateResult:
        """기술 게이트 자동 검증"""
        checks = {}
        for metric, value in performance_metrics.items():
            threshold = performance_thresholds.get(metric, 0)
            checks[f"성능_{metric}"] = value >= threshold
        checks["재현성"] = reproducibility_passed
        checks["지연시간"] = latency_ms <= latency_threshold
        all_passed = all(checks.values())
        result = GateResult(
            gate_name="기술 리뷰",
            status=(
                ApprovalStatus.APPROVED
                if all_passed
                else ApprovalStatus.REJECTED
            ),
            reviewer="자동 검증 시스템",
            timestamp=datetime.datetime.now().isoformat(),
            checks=checks,
        )
        self.gates.append(result)
        return result

    def run_fairness_gate(
        self,
        fairness_results: Dict[str, Dict],
    ) -> GateResult:
        """공정성 게이트 자동 검증"""
        checks = {}
        for attr, result in fairness_results.items():
            checks[f"공정성_{attr}_DP"] = result.get(
                "dp_pass", False
            )
            checks[f"공정성_{attr}_EO"] = result.get(
                "eo_pass", False
            )
        all_passed = all(checks.values())
        result = GateResult(
            gate_name="공정성 및 윤리 리뷰",
            status=(
                ApprovalStatus.APPROVED
                if all_passed
                else ApprovalStatus.REJECTED
            ),
            reviewer="공정성 검증 시스템",
            timestamp=datetime.datetime.now().isoformat(),
            checks=checks,
        )
        self.gates.append(result)
        return result

    def run_security_gate(
        self,
        robustness_passed: bool,
        vulnerability_scan_passed: bool,
        dependency_audit_passed: bool,
    ) -> GateResult:
        """보안 게이트 자동 검증"""
        checks = {
            "강건성_테스트": robustness_passed,
            "취약점_스캔": vulnerability_scan_passed,
            "의존성_감사": dependency_audit_passed,
        }
        all_passed = all(checks.values())
        result = GateResult(
            gate_name="보안 리뷰",
            status=(
                ApprovalStatus.APPROVED
                if all_passed
                else ApprovalStatus.REJECTED
            ),
            reviewer="보안 검증 시스템",
            timestamp=datetime.datetime.now().isoformat(),
            checks=checks,
        )
        self.gates.append(result)
        return result

    def add_manual_gate(
        self,
        gate_name: str,
        reviewer: str,
        status: ApprovalStatus,
        comments: str = "",
        conditions: List[str] = None,
    ) -> GateResult:
        """수동 승인 게이트 추가 (법무 등)"""
        result = GateResult(
            gate_name=gate_name,
            status=status,
            reviewer=reviewer,
            timestamp=datetime.datetime.now().isoformat(),
            comments=comments,
            conditions=conditions or [],
        )
        self.gates.append(result)
        return result

    def compute_final_decision(self) -> ApprovalStatus:
        """최종 배포 결정"""
        if not self.gates:
            return ApprovalStatus.PENDING
        statuses = [g.status for g in self.gates]
        if any(s == ApprovalStatus.REJECTED for s in statuses):
            self.final_status = ApprovalStatus.REJECTED
        elif any(
            s == ApprovalStatus.CONDITIONAL for s in statuses
        ):
            self.final_status = ApprovalStatus.CONDITIONAL
        elif all(
            s == ApprovalStatus.APPROVED for s in statuses
        ):
            self.final_status = ApprovalStatus.APPROVED
        else:
            self.final_status = ApprovalStatus.PENDING
        return self.final_status

    def generate_approval_report(
        self, output_path: str = None,
    ) -> Dict:
        """승인 보고서 생성"""
        self.compute_final_decision()
        report = {
            "model_name": self.model_name,
            "model_version": self.model_version,
            "final_decision": self.final_status.value,
            "timestamp": datetime.datetime.now().isoformat(),
            "gates": [],
        }
        for gate in self.gates:
            gate_info = {
                "name": gate.gate_name,
                "status": gate.status.value,
                "reviewer": gate.reviewer,
                "timestamp": gate.timestamp,
                "checks": gate.checks,
                "comments": gate.comments,
            }
            if gate.conditions:
                gate_info["conditions"] = gate.conditions
            report["gates"].append(gate_info)
        if output_path:
            with open(output_path, "w", encoding="utf-8") as f:
                json.dump(
                    report, f, ensure_ascii=False, indent=2
                )
        return report

# 사용 예시
system = DeploymentApprovalSystem(
    model_name="CreditScoring",
    model_version="2.1.0",
)
system.run_technical_gate(
    performance_metrics={"auc": 0.87, "f1": 0.82},
    performance_thresholds={"auc": 0.80, "f1": 0.75},
    reproducibility_passed=True,
    latency_ms=45.0,
    latency_threshold=100.0,
)
system.run_fairness_gate({
    "gender": {"dp_pass": True, "eo_pass": True},
    "age": {"dp_pass": True, "eo_pass": False},
})
system.run_security_gate(
    robustness_passed=True,
    vulnerability_scan_passed=True,
    dependency_audit_passed=True,
)
system.add_manual_gate(
    gate_name="법무 리뷰",
    reviewer="이법무 변호사",
    status=ApprovalStatus.APPROVED,
    comments="개인정보보호법 준수 확인 완료",
)
report = system.generate_approval_report(
    "approval_report_credit_v2.json"
)
decision = system.compute_final_decision()
print(f"최종 배포 결정: {decision.value}")
\end{lstlisting}

%----------------------------------------------------------
\section{기술 부채 관리}
\label{sec:10.7}
%----------------------------------------------------------

\cite{sculley2015hidden}가 경고했듯이, ML 시스템은 일반 소프트웨어보다 기술 부채(Technical Debt)가 훨씬 빠르게 축적된다. 실제 ML 코드는 전체 시스템의 극히 일부에 불과하며, 나머지는 데이터 수집, 특성 추출, 설정 관리, 모니터링 등의 ``배관 작업(Plumbing)''이 차지한다. 이 배관 작업에서 기술 부채가 주로 발생한다.

\subsection{ML 기술 부채의 유형}
\label{sec:10.7.1}

\paragraph{접착 코드(Glue Code)}
범용 ML 패키지를 실제 시스템에 통합하기 위해 작성하는 대량의 접착 코드는 유지보수의 악몽이 된다. 패키지 업데이트 시 호환성이 깨지기 쉽고, 테스트가 어렵다.

\paragraph{파이프라인 정글(Pipeline Jungle)}
데이터 전처리 파이프라인이 점점 복잡해지면서 스파게티 코드처럼 얽히는 현상이다. 새로운 데이터 소스가 추가될 때마다 임시 방편으로 코드가 덧붙여져 전체 구조를 파악하기 어려워진다.

\paragraph{숨겨진 피드백 루프(Hidden Feedback Loops)}
모델의 출력이 다시 입력 데이터에 영향을 미치는 순환 구조가 의도치 않게 형성되는 것이다. 예를 들어, 추천 시스템이 노출한 콘텐츠가 사용자 행동을 변화시키고, 이 변화된 행동 데이터가 다시 모델 학습에 사용된다.

\paragraph{미사용 특성(Undeclared Consumers)}
더 이상 유용하지 않은 특성이 제거되지 않고 남아 있거나, 어떤 시스템이 해당 특성에 의존하는지 파악되지 않는 상황이다.

\paragraph{설정 부채(Configuration Debt)}
하이퍼파라미터, 특성 선택, 데이터 경로 등의 설정이 체계적으로 관리되지 않아 ``어떤 설정이 현재 프로덕션에서 사용 중인지'' 알 수 없게 되는 상황이다.

\subsection{기술 부채 측정과 관리 전략}
\label{sec:10.7.2}

기술 부채는 측정하지 않으면 관리할 수 없다. 다음과 같은 지표를 정기적으로 측정한다.

\begin{itemize}
  \item \textbf{코드 복잡도}: 순환 복잡도(Cyclomatic Complexity), 코드 라인 수 대비 테스트 커버리지
  \item \textbf{의존성 건강도}: 오래된 의존성 패키지 비율, 알려진 취약점이 있는 패키지 수
  \item \textbf{파이프라인 복잡도}: 데이터 파이프라인의 노드 수와 엣지 수, DAG의 깊이
  \item \textbf{문서화 현행화율}: 코드 변경 대비 문서 업데이트 비율
  \item \textbf{재학습 안정성}: 재학습 파이프라인의 실패율, 평균 복구 시간
\end{itemize}

관리 전략은 다음과 같다.
\begin{itemize}
  \item \textbf{부채 스프린트}: 정기적으로 기술 부채 해소를 위한 전담 스프린트를 배정 (예: 분기별 1주)
  \item \textbf{부채 예산}: 새로운 기능 개발과 부채 해소의 비율을 사전에 합의 (예: 70:30)
  \item \textbf{자동화된 부채 탐지}: 코드 품질 도구(SonarQube, pylint 등)를 CI/CD에 통합하여 자동 탐지
  \item \textbf{아키텍처 리뷰}: 분기별 아키텍처 리뷰를 통해 구조적 부채를 식별하고 리팩토링 계획 수립
\end{itemize}

\subsection{기술 부채 유형별 영향과 해결 방안}
\label{sec:10.7.3}

\begin{table}[htbp]
\centering
\caption{ML 기술 부채 유형별 영향과 해결 방안}
\label{tab:tech_debt}
\begin{tabularx}{\textwidth}{p{2.5cm}p{3cm}X}
\toprule
\textbf{부채 유형} & \textbf{영향} & \textbf{해결 방안} \\
\midrule
접착 코드 & 유지보수 비용 증가, 패키지 업그레이드 어려움 & 래퍼(Wrapper) 패턴으로 추상화 계층 구축. 인터페이스 표준화 \\
\midrule
파이프라인 정글 & 디버깅 어려움, 신규 인력 온보딩 지연 & 파이프라인 오케스트레이터(Airflow, Kubeflow) 도입. DAG 구조 정리 \\
\midrule
숨겨진 피드백 루프 & 예측 불가능한 모델 행동, 편향 증폭 & 피드백 루프 매핑 및 문서화. 오프라인 평가와 온라인 평가 분리 \\
\midrule
미사용 특성 & 불필요한 계산 비용, 유지보수 부담 & 특성 사용 추적 시스템 구축. 정기적 특성 감사(Feature Audit) \\
\midrule
설정 부채 & 설정 오류로 인한 장애, 재현 불가 & 설정 관리 시스템(예: Hydra, OmegaConf) 도입. 설정 스키마 검증 자동화 \\
\midrule
테스트 부채 & 회귀 버그 빈발, 배포 신뢰도 저하 & ML 전용 테스트 프레임워크 도입. 데이터 테스트, 모델 테스트, 인프라 테스트 구분 \\
\midrule
문서화 부채 & 지식 이전 실패, 규제 감사 대응 어려움 & 코드-문서 동기화 도구. 모델 카드 자동 생성. 변경 로그 자동화 \\
\bottomrule
\end{tabularx}
\end{table}

%----------------------------------------------------------
\section{실무 도구}
\label{sec:10.8}
%----------------------------------------------------------

\begin{practicaltool}{모델 개발 거버넌스 종합 체크리스트}
다음 체크리스트는 모델 개발의 각 단계에서 거버넌스 요구사항을 빠짐없이 확인하기 위한 도구이다. 각 항목의 완료 여부를 체크하고, 미완료 항목에 대한 사유와 계획을 기록한다.

\begin{tcolorbox}[colback=white, colframe=black!75!white, title=모델 개발 거버넌스 체크리스트]

\textbf{1. 개발 준비 단계}
\begin{itemize}
  \item[$\square$] 모델 리스크 등급이 분류되었는가?
  \item[$\square$] 데이터시트가 작성되고 검토되었는가?\cite{gebru2021datasheets}
  \item[$\square$] 개발 환경이 표준 컨테이너 이미지 기반으로 설정되었는가?
  \item[$\square$] 의존성 패키지가 승인된 레지스트리에서 설치되었는가?
  \item[$\square$] 공정성 지표와 임계값이 사전 합의되었는가?
\end{itemize}

\textbf{2. 실험 및 개발 단계}
\begin{itemize}
  \item[$\square$] 실험 추적 시스템(MLflow 등)에 모든 실험이 기록되고 있는가?
  \item[$\square$] 랜덤 시드가 고정되고 기록되었는가?
  \item[$\square$] 하이퍼파라미터가 설정 파일로 관리되고 있는가?
  \item[$\square$] 데이터 버전이 DVC 등으로 추적되고 있는가?
  \item[$\square$] 코드 리뷰가 ML 전용 체크리스트에 따라 수행되었는가?
\end{itemize}

\textbf{3. 검증 단계}
\begin{itemize}
  \item[$\square$] 성능 지표가 KPI를 달성하였는가?
  \item[$\square$] 공정성 검증이 모든 민감 속성에 대해 수행되었는가?
  \item[$\square$] 적대적 강건성 테스트가 수행되었는가?\cite{goodfellow2014explaining}
  \item[$\square$] 설명가능성이 리스크 등급에 맞게 구현되었는가?\cite{arrieta2020explainable}
  \item[$\square$] 입력 검증 및 이상치 탐지 로직이 구현되었는가?
\end{itemize}

\textbf{4. 승인 및 배포 단계}
\begin{itemize}
  \item[$\square$] 모델 카드가 작성되고 최신 상태인가?\cite{mitchell2019model}
  \item[$\square$] 다단계 리뷰(기술, 윤리, 법무, 보안)를 모두 통과하였는가?
  \item[$\square$] 카나리 배포 계획이 수립되었는가?
  \item[$\square$] 롤백 절차가 문서화되고 테스트되었는가?
  \item[$\square$] 모니터링 대시보드가 구축되었는가?
\end{itemize}

\textbf{5. 지속 관리}
\begin{itemize}
  \item[$\square$] 기술 부채 현황이 파악되고 해소 계획이 있는가?\cite{sculley2015hidden}
  \item[$\square$] 문서화가 코드 변경과 동기화되고 있는가?
  \item[$\square$] 재학습 파이프라인이 자동화되어 있는가?
\end{itemize}
\end{tcolorbox}
\end{practicaltool}

\begin{practicaltool}{모델 카드 템플릿 (확장판)}
\begin{tcolorbox}[colback=white, colframe=black!75!white, title=Model Card for {[모델명]}]

\textbf{1. 모델 개요}
\begin{itemize}
  \item 모델명 / 버전: Credit-Scoring-Alpha / v2.1.0
  \item 개발자: 김이박 (리스크 모델링팀)
  \item 모델 날짜: 2026-01-20
  \item 모델 유형: XGBoost Classifier
  \item 프레임워크: XGBoost 1.7.6, Python 3.10
  \item 파라미터 수: 약 2,300개 리프 노드
  \item 라이선스: 내부 전용 (대외 배포 불가)
\end{itemize}

\textbf{2. 의도된 용도}
\begin{itemize}
  \item 주 용도: 소상공인 대상 신용대출 심사 보조 (최종 결정은 심사역이 수행)
  \item 대상 사용자: 여신 심사팀 심사역
  \item 제한 사항: 개인(가계) 대출 심사에는 사용 불가. 자동 결정(무심사) 용도로 사용 불가
\end{itemize}

\textbf{3. 학습 데이터}
\begin{itemize}
  \item 데이터셋: 2024-2025 소상공인 대출 이력 데이터
  \item 데이터 크기: 학습 80,000건 / 검증 10,000건 / 테스트 10,000건
  \item 데이터 출처: 내부 여신 시스템 (데이터시트 ID: DS-2025-0142)
  \item 전처리: 결측값 중위수 대체, 범주형 원핫 인코딩, 표준화
  \item 데이터 버전: DVC 해시 a3f7b2c
\end{itemize}

\textbf{4. 성능 평가}
\begin{itemize}
  \item 평가 데이터셋: 2025 하반기 대출 이력 (테스트 셋 10,000건)
  \item 주요 성능 지표:
  \begin{itemize}
    \item AUC: 0.87 (목표: 0.80 이상)
    \item KS 통계량: 0.48 (목표: 0.40 이상)
    \item F1 점수: 0.82 (목표: 0.75 이상)
  \end{itemize}
  \item 신뢰 구간: AUC 95\% CI [0.85, 0.89]
\end{itemize}

\textbf{5. 공정성 평가}
\begin{itemize}
  \item 성별: 승인율 차이 2.8\%p (허용범위 5\%p 이내 충족)
  \item 연령대: 20대 대비 60대 승인율 비 0.88 (4/5 규칙 충족)
  \item 업종: 요식업 그룹의 정밀도 저하 확인 (후속 조치 필요)
  \item 검증 도구: FairnessValidationPipeline v1.2
\end{itemize}

\textbf{6. 설명가능성}
\begin{itemize}
  \item 글로벌: SHAP Summary Plot 제공, 상위 10개 특성 중요도 공개
  \item 로컬: 개별 심사 건에 대한 SHAP Waterfall Plot 제공
  \item 반사실적 설명: ``어떤 조건이 달라졌다면 결과가 바뀌는지'' 안내 기능 제공
\end{itemize}

\textbf{7. 강건성 및 보안}
\begin{itemize}
  \item 가우시안 노이즈 테스트: 통과 (강건성 점수 0.94)
  \item 입력 범위 검증: 구현 완료 (범위 외 입력 자동 거부)
  \item 의존성 취약점 스캔: 0건 (2026-01-18 기준)
\end{itemize}

\textbf{8. 윤리적 고려 및 제약사항}
\begin{itemize}
  \item 알려진 편향: 특정 업종(요식업)에 대한 데이터 부족으로 예측 불안정 가능성
  \item 완화 조치: 요식업종 신청 건은 심사역 수동 리뷰 필수화
  \item 모니터링 계획: 월 1회 공정성 지표 모니터링, 분기별 전체 재검증
\end{itemize}

\textbf{9. 승인 이력}
\begin{itemize}
  \item 기술 리뷰: 2026-01-22 통과 (리뷰어: 박데이터)
  \item 공정성 리뷰: 2026-01-23 조건부 통과 (요식업 모니터링 조건)
  \item 법무 리뷰: 2026-01-24 통과 (리뷰어: 이법무)
  \item 보안 리뷰: 2026-01-24 통과 (리뷰어: 정보안)
  \item 최종 승인: 2026-01-25 조건부 승인 (카나리 배포 10\%로 시작)
\end{itemize}
\end{tcolorbox}
\end{practicaltool}

%----------------------------------------------------------
\subsection*{제10장 요약}
%----------------------------------------------------------

모델 개발 단계의 거버넌스는 창의적인 모델 개발 활동과 책임 있는 AI 구현 사이의 균형을 잡는 것이 핵심이다. 본 장에서 다룬 내용을 요약하면 다음과 같다.

\begin{enumerate}
  \item \textbf{표준화된 문서화}: 모델 카드\cite{mitchell2019model}와 데이터시트\cite{gebru2021datasheets}를 자동화된 체계로 작성하여 투명성을 확보하고, 개발 환경을 컨테이너 기반으로 표준화하여 재현성의 기반을 마련한다.

  \item \textbf{실험 관리와 재현성}: MLflow, W\&B 등의 실험 추적 도구를 활용하여 모든 실험을 기록하고, 랜덤 시드 고정, 환경 스냅샷, 데이터 버전 관리를 통해 재현성을 보장한다.

  \item \textbf{공정성 검증}: 인구통계학적 동등성, 균등 기회, 보정, 개인 공정성 등 다양한 지표를 상황에 맞게 적용하고, 전처리/중간처리/후처리 기법으로 편향을 완화하며, 공정성-정확도 트레이드오프를 투명하게 관리한다.

  \item \textbf{강건성과 보안}: 적대적 공격(FGSM, PGD 등)\cite{goodfellow2014explaining}에 대한 방어 체계를 구축하고, 입력 검증, 이상치 탐지, 모델 추출 공격 방어 등 다층적 보안 체계를 갖춘다.

  \item \textbf{설명가능성 구현}: 리스크 등급에 따라 글로벌/로컬 설명가능성 수준을 차등 적용하고\cite{arrieta2020explainable}, SHAP, LIME, 반사실적 설명 등을 통합한 보고서를 자동 생성한다.

  \item \textbf{승인과 릴리스 관리}: 기술, 윤리, 법무, 보안의 다단계 승인 게이트를 거치고, 카나리 배포와 A/B 테스트를 통한 점진적 배포를 수행하며, 비상 시 롤백 절차를 수립한다.

  \item \textbf{기술 부채 관리}: 접착 코드, 파이프라인 정글, 숨겨진 피드백 루프 등 ML 고유의 기술 부채\cite{sculley2015hidden}를 식별하고, 체계적으로 측정하며 관리한다.
\end{enumerate}

이러한 거버넌스 체계는 일회성 프로젝트가 아니라 조직의 문화로 내재화되어야 한다. 모델 개발자가 거버넌스를 ``추가적인 부담''이 아닌 ``품질 높은 모델을 만들기 위한 당연한 과정''으로 인식할 때, 비로소 책임 있는 AI가 실현된다.
